% Notation of the paper. Table 1 of Section 2 lists the symbols and where each is defined.
% Section files do not define macros of their own; anything not here is written out.

% --- numbers, sets, calculus
\newcommand{\R}{\mathbb{R}}
\newcommand{\N}{\mathbb{N}}
\newcommand{\Z}{\mathbb{Z}}
\newcommand{\dd}{\mathrm{d}}
\newcommand{\abs}[1]{\lvert #1\rvert}
\newcommand{\norm}[1]{\lVert #1\rVert}
\newcommand{\ip}[2]{\langle #1,#2\rangle}          % trace form <X,Y> = tr(XY)
\DeclareMathOperator{\tr}{tr}
\DeclareMathOperator{\conv}{conv}
\DeclareMathOperator{\inter}{int}
\DeclareMathOperator{\area}{area}
\DeclareMathOperator{\vol}{vol}

% --- packings (Section 2.1)
\newcommand{\doct}{\delta_{\mathrm{oct}}}          % the octagon's density
\newcommand{\Aoct}{A_{\mathrm{oct}}}               % sqrt(12) * delta_oct, the octagon's cost
\newcommand{\Koct}{K_{\mathrm{oct}}}               % Reinhardt's smoothed octagon (the Challenge's)
\newcommand{\KH}{K_{\mathrm{H}}}                   % the octagon of Hales's trajectory

% --- the control problem (Section 2.3)
\newcommand{\SLtwo}{\mathrm{SL}_2(\R)}
\newcommand{\sltwo}{\mathfrak{sl}_2(\R)}
\newcommand{\hstar}{\mathfrak{h}^{\star}}          % star domain in the upper half-plane
\newcommand{\hstarstar}{\mathfrak{h}^{\star\star}} % Hales--Vajjha's compact domain
\newcommand{\UT}{U_{T}}                            % the control simplex
\newcommand{\ee}{\mathbf{e}}                       % vertices e_0, e_1, e_2 of the simplex
\newcommand{\cost}{\mathcal{C}}                    % the cost C(T, X)
\newcommand{\tf}{t_{\mathrm{f}}}                   % the octagon's terminal time
\newcommand{\tsw}{t_{\mathrm{sw}}}                 % the octagon's switching time

% --- periods (Section 4)
\newcommand{\cls}[1]{\ensuremath{\mathrm{#1}}}     % the classes R, F, M1, M2, M3

% --- Lean declarations
% \lean{name} typesets a declaration name (underscores allowed); \leanref{name} cites the formal
% counterpart of a statement, as in "(Lean: name)". Names are upright even in italic text, are not
% letter-spaced, and break only after a dot, an underscore, a hyphen or a slash.
\renewcommand{\UrlFont}{\ttfamily\upshape}
\Urlmuskip=0mu
\makeatletter
\def\UrlBreaks{\do\.\do\_\do\-\do\/}
\def\UrlBigBreaks{}
\def\UrlNoBreaks{}
\makeatother
\newcommand{\lean}[1]{\textup{\path{#1}}}
\newcommand{\leanref}[1]{\textup{(Lean:~\path{#1})}}

% --- run-in headings inside a section
\newcommand{\runin}[1]{\par\medskip\noindent\emph{#1}\enspace}
